% ==================== 第 1 章 绪论 ====================
\section{绪\quad 论}

\subsection{项目背景}

\subsubsection{行业背景：K12 一对一辅导的"排课难"}

在个性化教学场景中，课程安排需要同时协调教师资质、学生可用时间、
教室容量与设备等资源。与固定行政班课表相比，一对一及小班辅导的
需求更分散、变更更频繁，因此排课质量直接影响履约效率和服务体验。
一对一业务的典型特征是：

\begin{enumerate}[leftmargin=2em,itemsep=0.2em]
\item \textbf{师资不稳定：}优秀教师稀缺、跨校区共享、排班灵活；
\item \textbf{教室资源紧张：}同一时段多班级争抢有限教室；
\item \textbf{学生时间碎片化：}晚自习、艺考、竞赛导致学生可上课时间
高度分散；
\item \textbf{动态变化频繁：}请假、补课、转班、调休每天都在发生。
\end{enumerate}

传统流程通常依赖电子表格和即时通信工具反复核对。一旦业务规模扩大，
人工方式容易出现资源冲突、课次遗漏和变更链路不透明等问题。本项目
因此把“零硬冲突”作为可验证的底线，而不使用未经留档的行业比例支撑结论。

\subsubsection{技术背景：从贪心到约束规划}

排课问题（Timetabling Problem）自 20 世纪 60 年代起就是
运筹学与人工智能领域的经典问题。1995 年 Werra 将其形式化为
\textbf{图着色问题}；2002 年 Burke 等人系统综述了
遗传算法、模拟退火、禁忌搜索等元启发式方法；
Google OR-Tools 提供的 \textbf{CP-SAT} 求解器将约束传播、SAT 搜索与
整数规划技术结合，为离散排课问题提供了成熟的工程实现。

与贪心/遗传算法相比，CP-SAT 在排课问题上的核心优势是
\textbf{硬约束严格保证}：教师/教室/班级的"同一时段只能排一节课"
不再是"尽量满足"而是"必须满足"。这对于教务场景至关重要，
因为任何冲突都会直接导致线下教学事故。

\subsubsection{飞书生态背景：低代码 + 多维表格 + 消息卡片}

飞书（Feishu/Lark）作为新一代企业协作平台，已在教育行业
积累了多维表格（Base）、审批、消息卡片、AI 助手等成熟能力。
本项目选择"原生算法 + 飞书集成"的路径，把核心排课能力
做成可被飞书多维表格直接调用的"智能体能力"，教务老师在
飞书侧就能完成"导入数据 → 看到三套方案 → 一键发布 → 通知师生"
的完整闭环，无需切换系统。

\subsection{项目意义}

\subsubsection{理论意义}

本项目主要形成以下三方面方法性贡献：
\begin{enumerate}[leftmargin=2em,itemsep=0.2em]
\item \textbf{多目标排课的统一框架：}将偏好、容量、稳定性、变更数
六类软目标统一到 CP-SAT 线性目标函数中，给出可比较的
策略权重表，并通过可运行原型验证其工程可行性。
\item \textbf{稳定性目标的线性化技术：}教师日均课次方差本质是非线性的，
本项目以绝对偏差代理方差目标，并通过辅助变量进行线性化，
为类似"分组均衡"软目标提供通用解法。
\item \textbf{最小影响调课的求解器实现：}把"变更数"作为软目标权重最高的
维度，让 CP-SAT 自行最小化变更，验证约束规划处理局部重排问题的可行性。
\end{enumerate}

\subsubsection{实践意义}

\begin{enumerate}[leftmargin=2em,itemsep=0.2em]
\item \textbf{降本潜力：}减少人工枚举组合与重复核对，把教务人员的精力
转移到规则确认、异常处理和方案审批。
\item \textbf{质量保障：}对求解结果执行独立校验；当前演示数据的 8 个课次
均被完整安排，硬冲突数为 0。
\item \textbf{响应提速：}教师请假后自动生成新方案及变更明细，为后续
消息通知和审批流提供结构化输入。
\end{enumerate}

\subsection{国内外研究现状}

\subsubsection{国外研究}

国际排课研究起源于高校课表（University Course Timetabling），
代表数据集包括 ITC 2007、Curriculum-Based Timetabling。
Google OR-Tools 团队在 2018--2024 年间持续维护 CP-SAT 求解器，
并持续用于各类组合优化问题。商业产品方面，
德国 \textbf{Untis}、英国 \textbf{TimeTabler}、澳大利亚
\textbf{aSc Timetables} 已形成成熟市场，但其
"硬约束求解 + 多策略对比"能力相对薄弱，且与国内教务流程
脱节严重。

\subsubsection{国内研究}

国内研究主要分两支：一支以高校课表为对象（如清华、北大等
高校自主开发排课系统），另一支以中小学/课外辅导为对象
（如校宝、晓黑板、ClassIn 等 SaaS）。
前者侧重科研，对商业体验关注不足；后者侧重产品体验，
在硬约束求解上仍主要依赖规则引擎，在大班次/多策略对比上
存在明显瓶颈。

本项目定位于\textbf{两者之间的中间地带}：使用工业级求解器
（OR-Tools CP-SAT）保证硬约束，使用多策略权重表满足
差异化运营需求，使用飞书生态降低部署与传播门槛。

\subsection{本项目主要工作与贡献}

本项目在 6 周时间内完成了 MVP 第一阶段开发，主要工作包括：
\begin{enumerate}[leftmargin=2em,itemsep=0.2em]
\item 设计并实现基于 \textbf{OR-Tools CP-SAT} 的核心排课求解器，
支持 5 套候选策略；
\item 设计 \textbf{5 维度策略权重表}（学生/教师/效率/平衡/重排），
使同一份业务数据可生成多套可比较方案；
\item 提出 \textbf{教师日均方差软目标}及线性化方案；
\item 实现 \textbf{最小影响调课}算法与版本差异比较；
\item 实现 \textbf{REST API}（6 个端点）与 4 个独立子页面；
\item 编写 \textbf{19 个自动化测试}（12 个求解器测试、7 个 API 集成测试），
搭建 GitHub Actions 双版本 CI 流水线；
\item 编写 \textbf{演示数据、批量启动脚本、停止脚本}，支持
双击 \code{launch\_classmind.bat} 一键启动。
\end{enumerate}

\subsection{论文组织结构}

本报告共分 9 章，各章内容安排如下：
\begin{itemize}[leftmargin=2em,itemsep=0.1em]
\item \textbf{第 1 章 绪论}：介绍项目背景、意义、国内外现状与本文工作。
\item \textbf{第 2 章 需求分析与可行性研究}：业务需求、用户需求、功能需求、
非功能需求、技术/经济/社会可行性。
\item \textbf{第 3 章 关键技术综述}：约束规划、OR-Tools CP-SAT、
Python 数据类、BaseHTTPRequestHandler、JSON 数据格式等。
\item \textbf{第 4 章 总体设计}：系统架构、模块划分、技术选型、数据流。
\item \textbf{第 5 章 核心算法设计}：决策变量、硬约束、软目标、
策略权重表、稳定性线性化、调课算法。
\item \textbf{第 6 章 系统详细设计与实现}：求解器、API、UI、CI。
\item \textbf{第 7 章 测试与结果分析}：单元测试、集成测试、性能分析。
\item \textbf{第 8 章 项目管理}：进度计划、团队分工、风险与对策。
\item \textbf{第 9 章 总结与展望}：阶段成果、第二阶段规划。
\end{itemize}

报告最后附有附录（数据结构定义、接口列表、关键源码、
参考文献与致谢）。
